\documentclass[otchet]{SCWorks}
% Тип обучения (одно из значений):
%    bachelor   - бакалавриат (по умолчанию)
%    spec       - специальность
%    master     - магистратура
% Форма обучения (одно из значений):
%    och        - очное (по умолчанию)
%    zaoch      - заочное
% Тип работы (одно из значений):
%    coursework - курсовая работа (по умолчанию)
%    referat    - реферат
%  * otchet     - универсальный отчет
%  * nirjournal - журнал НИР
%  * digital    - итоговая работа для цифровой кафедры
%    diploma    - дипломная работа
%    pract      - отчет о научно-исследовательской работе
%    autoref    - автореферат выпускной работы
%    assignment - задание на выпускную квалификационную работу
%    review     - отзыв руководителя
%    critique   - рецензия на выпускную работу

% * Добавлены вручную. За вопросами к @mchernigin
\usepackage{preamble}

\begin{document}

% Кафедра (в родительном падеже)
\chair{математической кибернетики и компьютерных наук}

% Тема работы
\title{Анализ сложности быстрой и пирамидальной сортировок}

% Курс
\course{2}

% Группа
\group{211}

% Факультет (в родительном падеже) (по умолчанию "факультета КНиИТ")
% \department{факультета КНиИТ}

% Специальность/направление код - наименование
 \napravlenie{02.03.02 "--- Фундаментальная информатика и информационные технологии}
% \napravlenie{02.03.01 "--- Математическое обеспечение и администрирование информационных систем}
% \napravlenie{09.03.01 "--- Информатика и вычислительная техника}
%\napravlenie{09.03.04 "--- Программная инженерия}
% \napravlenie{10.05.01 "--- Компьютерная безопасность}

% Для студентки. Для работы студента следующая команда не нужна.
% \studenttitle{студентки}

% Фамилия, имя, отчество в родительном падеже
\author{Бурова Арсения Александровича}

% Заведующий кафедрой 
\chtitle{доцент, к.\,ф.-м.\,н.}
\chname{С.\,В.\,Миронов}

% Руководитель ДПП ПП для цифровой кафедры (перекрывает заведующего кафедры)
% \chpretitle{
%     заведующий кафедрой математических основ информатики и олимпиадного\\
%     программирования на базе МАОУ <<Ф"=Т лицей №1>>
% }
% \chtitle{г. Саратов, к.\,ф.-м.\,н., доцент}
% \chname{Кондратова\, Ю.\,Н.}

% Научный руководитель (для реферата преподаватель проверяющий работу)
\satitle{доцент, к.\,ф.-м.\,н.} %должность, степень, звание
\saname{Ю.\,Н.\,Кондратова}

% Руководитель практики от организации (руководитель для цифровой кафедры)
\patitle{доцент, к.\,ф.-м.\,н.}
\paname{С.\,В.\,Миронов}

% Руководитель НИР
\nirtitle{доцент, к.\,п.\,н.} % степень, звание
\nirname{В.\,А.\,Векслер}

% Семестр (только для практики, для остальных типов работ не используется)
\term{2}

% Наименование практики (только для практики, для остальных типов работ не
% используется)
\practtype{учебная}

% Продолжительность практики (количество недель) (только для практики, для
% остальных типов работ не используется)
\duration{2}

% Даты начала и окончания практики (только для практики, для остальных типов
% работ не используется)
\practStart{01.07.2024}
\practFinish{13.01.2024}

% Год выполнения отчета
\date{2026}

\maketitle

% Включение нумерации рисунков, формул и таблиц по разделам (по умолчанию -
% нумерация сквозная) (допускается оба вида нумерации)
\secNumbering

\tableofcontents

% Раздел "Обозначения и сокращения". Может отсутствовать в работе
% \abbreviations
% \begin{description}
%     \item ... "--- ...
%     \item ... "--- ...
% \end{description}

% Раздел "Определения". Может отсутствовать в работе
% \definitions

% Раздел "Определения, обозначения и сокращения". Может отсутствовать в работе.
% Если присутствует, то заменяет собой разделы "Обозначения и сокращения" и
% "Определения"
% \defabbr

\section{Реализация алгоритмов}
\subsection{Быстрая сортировка}
\begin{Verbatim}[
  numbers=left,
  breaklines=true,
  breakanywhere=true
]
void quickSort(int arr[], int L, int R) {
    int i = L; //начальные индексы
    int j = R;
    int op = arr[(L + R) / 2]; //выбираем опорный элемент из середины
    while (i <= j) { //основной цикл разделения
        //двигаем i вправо, пока элементы меньше опорного
        while (arr[i] < op) {
            i++;
        }
        //двигаем j влево, пока элементы больше опорного
        while (arr[j] > op) {
            j--;
        }
        //если индексы не пересеклись, меняем элементы
        if (i <= j) {
            swap(arr[i], arr[j]);
            i++;
            j--;
        }
    }
    //рекурсивно сортируем левую часть
    if (L < j) {
        quickSort(arr, L, j);
    }
    //рекурсивно сортируем правую часть
    if (i < R) {
        quickSort(arr, i, R);
    }
}
}
\end{Verbatim}

\subsection{Пирамидальная сортировка}

\begin{Verbatim}[
  numbers=left,
  breaklines=true,
  breakanywhere=true
]
void heapify(int arr[], int i, int N) {
    //полагаем, что i-тый элемент является максимальным
    int parent = i;
    while (true) {
        int leftChild = 2 * i + 1; //левый ребенок
        int rightChild = 2 * i + 2; //правый ребенок
        int maxElement = parent;

        //если существует левый ребенок и он больше текущего максимума
        if (leftChild < N && arr[leftChild] > arr[maxElement]) {
            maxElement = leftChild;
        }

        //если существует правый ребенок и он больше текущего максимума
        if (rightChild < N && arr[rightChild] > arr[maxElement]) {
            maxElement = rightChild;
        }

        //если максимум - это родитель, прекращаем цикл
        if (maxElement == parent) {
            break;
        }

        //иначе меняем местами максимум и родителя
        swap(arr[parent], arr[maxElement]);

        //продолжаем просеивание с позиции maxElement
        i = maxElement;
        parent = i;
    }
}
\end{Verbatim}

\begin{Verbatim}[
  numbers=left,
  breaklines=true,
  breakanywhere=true
]
void heapSort(int arr[], int N) {
    for (int i = N / 2 - 1; i >= 0; i--) {
        heapify(arr, i, N);
    }
    for (int i = N - 1; i >= 1; i--) {
        //меняем местами корень с последним
        swap(arr[0], arr[i]);

        heapify(arr, 0, i);
    }
}
\end{Verbatim}


\newpage

\section{Анализ сложности алгоритмов}

\subsection{Быстрая сортировка}

Для двух вложенных циклов \texttt{while (arr[i] < op)} и \texttt{while (arr[j] > op)} общее количество перемещений указателей $i$ и $j$ в сумме не превышает длины текущего подмассива, то есть $(t + m) = O(n)$, где $n = R - L + 1$.

Внешний цикл \texttt{while (i <= j)} выполняется $k$ раз, но на каждой итерации хотя бы один из указателей $i$ или $j$ сдвигается. Следовательно, общее число итераций внешнего цикла также ограничено $O(n)$. Таким образом, трудоёмкость этапа разделения составляет $O(n)$.

Пусть $T(n)$ --- время работы алгоритма для массива из $n$ элементов. После разделения левая часть содержит $j - L + 1$ элементов, правая — $R - i + 1$ элементов. В лучшем случае эти части примерно равны.

Построим рекуррентное соотношение для лучшего случая:
\[
T(n) = 
\begin{cases} 
c, & \text{если } n = 1, \\
2T\left(\dfrac{n}{2}\right) + bn, & \text{если } n > 1.
\end{cases}
\]

Здесь $a = 2$ --- количество подзадач, $n/k$ --- размер подзадачи ($k = 2$ — постоянная величина). Раскрываем:

\[
\begin{aligned}
T(n) &= 2T\left(\frac{n}{2}\right) + bn \\
     &= 2\left[2T\left(\frac{n}{4}\right) + b\frac{n}{2}\right] + bn = 4T\left(\frac{n}{4}\right) + bn + bn \\
     &= 4T\left(\frac{n}{4}\right) + 2bn \\
     &= \dots \\
     &= 2^k T\left(\frac{n}{2^k}\right) + k \cdot bn.
\end{aligned}
\]
При $k = \log_2 n$ получим $T(n) = n \cdot T(1) + bn \log_2 n = O(n \log n)$.

В худшем случае опорный элемент каждый раз оказывается самым большим или самым маленьким элементом подмассива, что приводит к неравномерному разделению. Один из подмассивов оказывается пустым, а другой содержит $n-1$ элемент.

Глубина рекурсии в таком случае становится равной $n$. На каждом уровне рекурсии размер обрабатываемого подмассива уменьшается на единицу: $n, n-1, n-2, \dots, 1$.

Суммарное количество операций:
\[
T(n) = n + (n-1) + (n-2) + \dots + 1 = \frac{n(n+1)}{2} = O(n^2).
\]
Сумму посчитали через формулу арифметической прогрессии.

Таким образом, в худшем случае трудоёмкость быстрой сортировки составляет $O(n^2)$.

\subsection{Пирамидальная сортировка}

Сортировка основана на построении кучи. Свойства кучи:
\begin{enumerate}
    \item Узел всегда больше своих потомков;
    \item На последнем уровне элементы располагаются слева направо, пока они не закончатся.
\end{enumerate}

Т.к. куча — это бинарное дерево, то его высота не больше, чем $\log_2 n$, а значит временная сложность функции просеивания \texttt{heapify} равна $O(\log n)$.

В первом цикле сортировки мы проходимся по всем нелистовым узлам (от $N/2 - 1$ до $0$), вызывая для каждого \texttt{heapify}. Количество таких узлов составляет $O(n)$, каждый вызов — $O(\log n)$. Следовательно, построение начальной кучи выполняется за:
\[
T = O(n \log n)
\]

Во втором цикле мы извлекаем элементы из кучи по одному (всего $N-1$ итераций), каждый раз восстанавливая свойства кучи вызовом \texttt{heapify}. Каждое восстановление выполняется за $O(\log n)$.

Общая сложность алгоритма:
\[
T(n) = O(n) + O(n \log n) = O(n \log n)
\]

Таким образом, пирамидальная сортировка в лучшем, среднем и худшем случаях имеет сложность:
\[
T(n) = O(n \log n)
\]

\end{document}
